\section{Related Work}
\label{sec:related}
Prior work on adaptive test-time compute spans diverse
mechanisms but rests on a shared, almost always implicit,
assumption: that uncertainty signals have a fixed,
setting-independent relationship with compute value. We organize
the literature around how each line of work commits to this
\emph{fixed-direction} assumption, since it is precisely what
this paper tests.

\paragraph{Gating on uncertainty signals.}
Signal-based methods allocate extra computation when token
confidence, entropy, or disagreement exceeds a threshold,
encoding the fixed-direction assumption directly in their
trigger rule.
In reasoning settings, this takes the form of extended
chain-of-thought generation~\cite{seag2025, cats2025,
corefine2026, thinkjustenough2025, s1_2025, han2025tokenbudget}.
In agent settings, vote-based methods~\cite{catts2026} invoke
an arbiter when candidate actions disagree, and verbalized
uncertainty methods~\cite{auq2026} trigger reflection when
self-reported confidence is low.

\paragraph{Gating on learned difficulty.}
Hidden-state probes~\cite{diffadapt2025, zhu2025llmalreadyknows}
estimate task difficulty from internal representations, and
RL-based methods~\cite{adaptthink2025, thinkless2025,
aggarwal2025l1, paglieri2025learning} learn think/no-think
policies end-to-end. 
% Though these methods do not commit to a
% specific signal, their gating policies effectively assume that
% higher difficulty warrants more compute, sharing the same fixed
% direction as signal-based approaches. 
% This assumption holds on
% homogeneous reasoning benchmarks (MATH, GSM8K, GPQA) where signal
% direction is consistent by construction. Whether it transfers to
% heterogeneous agent environments has not been directly tested; we
% show it does not (Proposition~\ref{prop:necessity}).
Although these methods commit to neither a specific signal nor an explicit threshold, their gating policies share the assumption that higher difficulty warrants more compute---identical in \emph{direction} to signal-based approaches. The combined body of work therefore rests on the same implicit fixed-direction assumption, and is evaluated almost exclusively on homogeneous reasoning benchmarks (GSM8K~\cite{DBLP:journals/corr/abs-2110-14168}, GPQA~\cite{DBLP:journals/corr/abs-2311-12022}) where the assumption holds by construction. We construct this assumption explicitly and test it in the heterogeneous agent setting, where it fails.

% \paragraph{Metareasoning and uncertainty semantics.}
% Rational metareasoning~\cite{russell1991right,
% desabbata2024metareasoning} formalizes the value of computation
% (VOC), classically assuming $\text{VOC} \geq 0$ since an agent can
% disregard unfavorable results. We show that this assumption fails
% in the agent setting, where the same model both proposes and
% evaluates actions: under this evaluator-executor identity, VOC can
% be \emph{negative}, making direction discovery necessary.
% \citet{tao2025revisiting} and \citet{heo2025llmuncertainty} find
% that calibration quality varies across task types, and
% \citet{yadkori2024believe} decouple epistemic and aleatoric
% uncertainty, providing converging evidence that uncertainty
% semantics are domain-dependent, though none documents direction
% reversal. Our two-source model formalizes this domain dependence
% (\S\ref{sec:toy-model}), and our reversal results make it
% actionable. Extended discussion is in Appendix~\ref{app:related}.
\paragraph{Metareasoning and uncertainty semantics.}
Rational metareasoning~\cite{russell1991right, desabbata2024metareasoning} formalizes the value of computation (VOC), classically assuming $\text{VOC} \geq 0$ since an agent can disregard unfavorable results. We show this assumption fails when the same model both proposes and evaluates actions (\S\ref{sec:toy-model}). \citet{tao2025revisiting} and \citet{heo2025llmuncertainty} find that calibration quality varies across task types, and \citet{yadkori2024believe} decouple epistemic and aleatoric uncertainty, providing converging evidence that uncertainty semantics are domain-dependent. 
However, none observes that signal direction can reverse across environments. Our two-source model formalizes this domain dependence and makes it actionable. Extended related work discussion is in Appendix~\ref{app:related}.
